\documentclass[12pt,a4paper]{article}
\usepackage[utf8]{inputenc}
\usepackage{amsmath,amssymb,amsthm}
\usepackage{algorithm}
\usepackage{algorithmic}
\usepackage{geometry}
\geometry{margin=1in}

\newtheorem{theorem}{Theorem}
\newtheorem{lemma}[theorem]{Lemma}
\newtheorem{proposition}[theorem]{Proposition}
\theoremstyle{definition}
\newtheorem{definition}[theorem]{Definition}

\title{Problem 28: Number Spiral Diagonals}
\author{}
\date{}

\begin{document}
\maketitle

\section{Problem Statement}

Starting with 1 at the center and spiraling clockwise, find the sum of the numbers on the diagonals in a $1001 \times 1001$ number spiral.

\section{Mathematical Development}

\begin{definition}[Spiral Rings]
\emph{Ring~$k$} ($k \geq 1$) is the layer whose outermost side has length $2k+1$. Ring~0 is the center cell containing~1.
\end{definition}

\begin{lemma}[Ring Boundaries]\label{lem:bounds}
Ring~$k$ contains the integers from $(2k-1)^2 + 1$ to $(2k+1)^2$. The count is $8k$.
\end{lemma}

\begin{proof}
$(2k+1)^2 - (2k-1)^2 = 8k$.
\end{proof}

\begin{theorem}[Corner Values]\label{thm:corners}
The four diagonal entries of ring~$k$ ($k \geq 1$) are
\begin{align*}
\mathrm{TR}(k) &= (2k+1)^2, & \mathrm{TL}(k) &= (2k+1)^2 - 2k, \\
\mathrm{BL}(k) &= (2k+1)^2 - 4k, & \mathrm{BR}(k) &= (2k+1)^2 - 6k.
\end{align*}
\end{theorem}

\begin{proof}
The last number in ring~$k$ is $(2k+1)^2$ at the top-right corner. Each side has $2k+1$ numbers but shares a corner with the next side, contributing $2k$ new numbers per side. Successive corners going counterclockwise are separated by~$2k$.
\end{proof}

\begin{theorem}[Ring Diagonal Sum]\label{thm:ringsum}
The sum of the four corners of ring~$k$ is $S_k = 16k^2 + 4k + 4$.
\end{theorem}

\begin{proof}
\[
S_k = \sum_{j=0}^{3}\bigl[(2k+1)^2 - 2jk\bigr] = 4(2k+1)^2 - 12k = 16k^2 + 4k + 4. \qedhere
\]
\end{proof}

\begin{theorem}[Closed-Form Formula]\label{thm:closed}
For an $N \times N$ spiral ($N$ odd), the diagonal sum is
\[
S(N) = \frac{4N^3 + 3N^2 + 8N - 9}{6}.
\]
\end{theorem}

\begin{proof}
Let $n = (N-1)/2$. Then
\begin{align*}
S &= 1 + \sum_{k=1}^{n}(16k^2 + 4k + 4) \\
  &= 1 + \frac{8n(n+1)(2n+1)}{3} + 2n(n+1) + 4n.
\end{align*}
Substituting $n = (N-1)/2$, $n+1 = (N+1)/2$, $2n+1 = N$:
\begin{align*}
S &= 1 + \frac{2N(N^2-1)}{3} + \frac{N^2-1}{2} + 2(N-1) \\
  &= \frac{6 + 4N^3 - 4N + 3N^2 - 3 + 12N - 12}{6} = \frac{4N^3 + 3N^2 + 8N - 9}{6}. \qedhere
\end{align*}
\end{proof}

\begin{lemma}[Verification]
$S(1) = 1$, $S(3) = 25$, $S(5) = 101$, matching direct computation.
\end{lemma}

\section{Algorithm}

\begin{algorithm}[H]
\caption{Spiral Diagonal Sum}
\begin{algorithmic}[1]
\REQUIRE Odd side length $N$
\ENSURE Sum of diagonal entries
\RETURN $(4N^3 + 3N^2 + 8N - 9) / 6$
\end{algorithmic}
\end{algorithm}

\section{Complexity Analysis}

\begin{proposition}
The closed-form evaluation runs in $O(1)$ time and $O(1)$ space.
\end{proposition}

\section{Answer}

\[
\boxed{669171001}
\]

\end{document}
